\documentclass{article}
\usepackage{amsmath,amssymb,hyperref}
\title{Editing as Control: Online Injection Scheduling for Causal Streaming Video Editors}
\author{}
\date{}
\begin{document}
\maketitle

\section{Motivation and problem statement}

We edit video with a \emph{frozen, pre-trained causal streaming generator} (the StreamEdit
setting): a few-step generator that emits frames left-to-right, each conditioned on the
realized past through a rolling KV cache. Editing is performed by a dual-branch scheme in
which source-video information is injected into the target branch through a self-attention
bridge (plus cross-attention grounding/boosting).

\textbf{The gap.} The injection strength is currently an \emph{open-loop} schedule: a fixed
function of diffusion timestep (and, optionally, a spatial mask), committed before generation
begins. A non-causal two-branch editor \emph{must} be open-loop, because all frames are
denoised jointly and there is no realized past to react to. Empirically, no single schedule is
best: the injection curve that maximizes quality depends on the clip and on the edit type
(local structure-preserving edit vs.\ global appearance/style vs.\ motion change). A fixed
schedule therefore leaves quality on the table, and hand-selecting a schedule per edit type
does not scale.

\textbf{What causality buys us.} After emitting frame $t$, a causal generator exposes a
realized state --- the decoded frame and the KV cache --- that can be \emph{measured}, so the
injection strength for frame $t{+}1$ can be a \emph{function of that measurement}. This closes
a feedback loop that is cheap precisely because generation is causal and few-step. The claim
is not that only causal models can obtain feedback (a non-causal editor could iterate over
multiple passes) but that causality makes \emph{single-pass, online, state-dependent} control
affordable.

\textbf{Problem.} Given a frozen streaming editor, a source video, and an edit instruction,
produce the injection control signal $u_t$ (per-frame, and ideally per-region and per-layer)
that maximizes edit quality --- jointly: source preservation in non-edited regions, edit
fidelity in edited regions, and temporal consistency --- \emph{online}, within the latency
budget that motivates streaming in the first place.

\section{Method}

We treat injection scheduling as \textbf{sequential decision-making over the injection control
variable}, conditioned on causal state. The paper has two coupled contributions, deliberately
ordered.

\subsection*{(1) The phenomenon: a training-free online rule (analysis + lower bound)}

A cheap, hand-designed feedback law that already beats every fixed schedule. At each streaming
step, compute a low-cost observation $o_t$ from the frozen model's own signals:
\begin{itemize}
  \item source-preservation error in the non-edited region (e.g.\ feature/statistic distance
        between source and target branch outside the edit mask);
  \item edit-prompt alignment in the edited region (e.g.\ cross-attention energy on the edit
        tokens, or a light CLIP-style score);
  \item temporal drift vs.\ the previous edited frame.
\end{itemize}
A fixed control law $u_{t+1} = g(o_t)$ (PID-like, or a small monotone rule) maps the
observation to the next injection strength. This is \emph{training-free and plug-and-play};
its role is to establish that the closed-loop gap is real and not a training artifact, and to
serve as an ablation baseline. \textbf{(speculative)} the exact observation set is unvalidated;
the design constraint is that computing $o_t$ must be cheap relative to a generation step.

\textbf{Where the proof lands.} Beating the dataset-averaged \emph{fixed-best} schedule is
necessary but \emph{not sufficient}: two gaps hide in ``online beats fixed.'' (i) A
\emph{per-clip adaptation} gap --- the schedule should depend on the clip --- is already
captured by cheaply \emph{selecting} the best fixed curve per clip, and has nothing to do with
causality. (ii) The \emph{within-clip closed-loop} gap is what a per-frame reactive controller
gets \emph{beyond any single fixed curve}, and only this requires state feedback. Since the
oracle per-clip schedule is still one fixed curve replayed on every frame, the rule's proof of
the closed-loop gap lands \emph{only against oracle-per-clip}, not against fixed-best: a
single-pass rule that beats the hindsight-chosen best fixed curve is clean evidence that
reacting to realized state buys something no fixed curve can. Beating only fixed-best collapses
the thesis into ``tune the schedule per clip,'' which needs no causality.

\subsection*{(2) The method: a learned injection-control policy}

A lightweight policy $\pi_\theta$ that reads causal state and outputs $u_t$ per frame (and per
region/layer), realizing the full potential the rule only approximates. Formulation:
\begin{itemize}
  \item \textbf{State.} The controller's observation is drawn from the streaming-specific
        state --- the KV cache summary and the decoded-so-far frames --- not just the diffusion
        timestep. \emph{This is the load-bearing use of causality and must remain so.}
  \item \textbf{Action.} $u_t$: injection gain(s), resolved per region and per attention layer.
  \item \textbf{Reward.} A trade-off of source preservation, edit fidelity, and temporal
        consistency, evaluated on the realized stream. \textbf{(speculative)} terminal vs.\
        dense reward shaping is unresolved; sparse terminal reward is the honest default but
        credit assignment over a long stream is a risk. The flag is about \emph{when} the policy
        receives reward during a rollout:
        \begin{itemize}
          \item \textbf{Terminal (sparse).} Score the edit \emph{once}, at clip end (final
                CLIP + preservation + temporal); one scalar per rollout. \emph{Pro:} the honest
                objective --- exactly what we care about, no per-frame proxy to game.
                \emph{Con:} credit assignment --- a single terminal scalar spread over $\sim$40
                per-frame $u_t$ decisions gives high-variance gradients and slow/unstable
                training (classic sparse-reward RL, worsened by the long horizon).
          \item \textbf{Dense (shaped).} Hand out a per-frame $r_t$ (e.g.\ this frame's
                preservation error + edit alignment). \emph{Pro:} easier credit assignment,
                faster convergence. \emph{Con:} optimizes a \emph{proxy} that may not sum to
                true perceptual quality --- a policy can maximize the shaped reward while
                degrading the edit (reward hacking; cf.\ ``Reward reliability'' below).
        \end{itemize}
        Default: sparse terminal reward (intellectually honest, matches GINH's setup), while
        acknowledging the long-horizon credit-assignment risk is real and may force some shaping.
  \item \textbf{Optional lookahead (MPC-lite).} Because the model is few-step, a cheap 1-step
        rollout under a few candidate settings lets the policy pick $u_t$ minimizing predicted
        next-frame error. \textbf{(speculative)} whether the lookahead cost is worth it under
        the latency budget is an open experiment. The flag is about whether to add a
        model-predictive-control step on top of (or instead of) the reactive policy:
        \begin{itemize}
          \item \textbf{The idea.} Since the generator is few-step, simulating one frame ahead
                is cheap. For frame $t$, try a handful of candidate injection values
                $\{u^{(1)}, u^{(2)}, \dots\}$, do a cheap 1-step rollout for each, measure
                predicted next-frame error, and pick the best --- planning rather than pure
                reaction. ``MPC-lite'' = a shallow (1-step, few-candidate) model-predictive
                control.
          \item \textbf{Spectrum point.} It sits \emph{between} the training-free rule and the
                learned policy, and can itself be \emph{training-free} (no $\pi_\theta$, just
                try-and-pick) --- a third rung on the fallback ladder (rule $\to$ MPC-lite $\to$
                imitation $\to$ RL).
          \item \textbf{Unresolved --- latency tension.} $K$ candidates means $K$ extra partial
                rollouts per frame; if that approaches the cost of a real generation step it
                eats the streaming latency budget --- the entire justification for a streaming
                editor. The quality gain may be real, but whether it survives the latency
                constraint is an open experiment (cf.\ ``Latency tension (existential)'' below),
                not a settled design choice.
        \end{itemize}
\end{itemize}
The generator stays frozen throughout; only $\pi_\theta$ is trained.

\subsection*{Design components (shared by both)}
\begin{itemize}
  \item \textbf{Per-region / per-layer resolution} of the injection gain, rather than a single
        global scalar per frame.
  \item \textbf{Drift-aware re-anchoring:} monitor divergence from the source's temporal
        structure and transiently boost source injection when drift is detected --- the
        test-time, injection-knob analog of trained drift correction.
\end{itemize}

\section{Novelty claims}

\begin{enumerate}
  \item \textbf{Closed-loop beats open-loop.} An online, state-dependent injection schedule
        strictly outperforms the best fixed schedule on a frozen streaming editor, across edit
        types. \emph{Falsifiable:} beating dataset-averaged fixed-best is only necessary; the
        real bar is \emph{oracle-per-clip} (the hindsight-chosen best fixed curve per clip). If
        online $\approx$ oracle-per-clip, the gain is mere per-clip tuning --- capturable by
        cheap selection, needing no causality --- and this claim fails.
  \item \textbf{Learned approaches oracle.} A learned policy conditioned on causal state
        approaches the oracle per-clip schedule (the upper bound), which the training-free rule
        does not fully reach. \emph{Falsifiable:} if the learned policy $\approx$ the
        training-free rule, the learned contribution is not justified.
  \item \textbf{Causality is load-bearing, not incidental.} Conditioning the controller on
        streaming-specific state (KV cache / realized frames) and/or cheap causal lookahead
        yields gains that a schedule ignoring that state cannot. This is what converts our
        contribution from an \emph{efficiency} claim (causality makes single-pass feedback
        cheap --- true but reducible to ``a cheaper application of GINH'') into a
        \emph{method-novelty} claim (the specific causal state is \emph{necessary}, not merely
        convenient). \emph{Falsifiable:} if a controller that sees only the diffusion timestep
        (GINH-shaped conditioning) matches the state-conditioned one, the streaming framing adds
        nothing over generic adaptive-guidance work and we are GINH on a new knob. \emph{This
        ablation --- timestep-only $<$ state-conditioned --- is the paper.}
  \item \textbf{Training-free, plug-and-play at the method level.} The whole approach sits on a
        frozen generator with no distillation or generator fine-tuning --- distinguishing it
        from trained streaming editors.
\end{enumerate}

\section{Relation to prior work}

\textbf{Adaptive/feedback control of guidance (mechanism precedent, wrong domain).}
CFG-Ctrl frames a guidance-weight scheduler as a time-varying proportional feedback controller
(gain scheduling). Feedback Guidance (FBG) uses a state-dependent coefficient to self-regulate
guidance. Adaptive Diffusion Guidance via Stochastic Optimal Control (Azangulov et al., 2505.19367)
casts scheduling as control selecting strength from time/sample/class. ``Guidance Is Not a
Hyperparameter'' (2605.07701) learns the schedule via RL as sequential decision-making under
sparse terminal reward. \emph{Difference:} all control the \emph{CFG guidance scalar} in
image / non-streaming diffusion. We control the \emph{source-injection strength} on a
streaming editor's self-attention bridge, conditioned on \emph{causal state}. This distinction
is the whole novelty --- if it blurs, claim 3 falls (see Open questions).

\textbf{Why GINH is the threat, precisely.} Strip the domain and our headline method (2) \emph{is}
GINH's abstraction: a learned policy setting a guidance-like scalar, framed as sequential
decision-making under sparse terminal reward, conditioned on state, generator frozen. The
one-line reviewer reduction is ``GINH applied to a different knob (injection vs.\ CFG) in a
different modality (streaming video edit vs.\ image gen).'' Crucially, \textbf{``we condition on
causal state'' is not the escape}: GINH already conditions its policy on state (sample, class,
time), and ``use realized state so far to set the next control step'' is exactly its move ---
ours merely runs the autoregression over frames rather than denoising steps. Our genuinely
distinct contribution is therefore an \emph{efficiency} one --- causality makes single-pass,
online, state-dependent feedback \emph{cheap}, where a non-causal editor would need multi-pass
iteration to obtain the same signal. That efficiency claim only upgrades to a \emph{method-novelty}
claim if the causal state is shown \emph{necessary}: a policy seeing only the diffusion timestep
(GINH-shaped conditioning) must fail to match the state-conditioned one (claim 3). That ablation
is the paper.

\textbf{Causal / streaming video editing (direct competitors, all trained).}
LiveEdit (ECCV 2026): causal frame-by-frame editing via three-stage distillation from a
bidirectional teacher, plus an AR-oriented mask cache. JoyAI-Video-Edit: 16B AR editor with
Source-Anchored DMD and long-horizon AR distillation to preserve source fidelity and mitigate
drift. AnchorEdit (2606.11751): AR-diffusion multi-turn \emph{image} editing with self-rollout
and anchor memory. \emph{Difference:} these train/distill the editor; we place a controller on
a \emph{frozen} generator and act at test time. Their existence makes the field crowded ---
our defensible niche is training-free test-time control.

\textbf{Drift / exposure bias in AR video (scaffolding for re-anchoring).}
Self Forcing (NeurIPS 2025) conditions each frame on self-generated outputs to fix exposure
bias; TetherCache is a training-free cache-management strategy that aligns recalled memory
tokens to a trusted distribution to reduce drift; Resampling Forcing / Self-Forcing++ extend
the toolkit. \emph{Difference:} these address drift in \emph{generation} via training or cache
edits; we address editing-time drift at test time via the injection knob.

\section{Open questions and risks}

\begin{itemize}
  \item \textbf{Latency tension (existential).} The justification for streaming is low latency.
        If computing $o_t$ / the reward / the lookahead costs near a generation step, the
        advantage is eaten. Controller and measurements must be cheap; this is a design
        constraint, not a footnote.
  \item \textbf{The oracle gap (make-or-break).} Required experiment: fixed-best schedule vs.\
        oracle per-clip schedule (upper bound) vs.\ training-free rule vs.\ learned policy. The
        load-bearing comparison is \emph{rule/policy vs.\ oracle-per-clip}, not vs.\ fixed-best:
        beating fixed-best only proves the schedule should be clip-dependent, which cheap
        best-of-$K$ selection already achieves without any feedback loop. To isolate the
        \emph{within-clip} closed-loop gap, the rule must beat oracle-per-clip single-pass (or
        an ablation freezing the rule to its clip-average $u$ must degrade). If online
        $\approx$ oracle-per-clip, the thesis collapses into per-clip hyperparameter sweeping.
  \item \textbf{Novelty vs.\ RL-guidance work (the central threat).} Our headline method \emph{is}
        GINH's abstraction (state-conditioned RL policy over a guidance-like scalar, sparse
        terminal reward, frozen model); ``we condition on causal state'' is not an escape, since
        GINH already conditions on state. The only distinct contribution is \emph{cheap
        single-pass feedback} (an efficiency axis GINH does not touch), and it upgrades to
        method-novelty \emph{only} by proving the causal state necessary: timestep-only $<$
        state-conditioned (claim 3). Absent that ablation, we are GINH applied to video editing.
  \item \textbf{``Only causal'' overclaim.} Frame the causal advantage as efficiency/latency
        (single-pass feedback), never as impossibility for non-causal editors.
  \item \textbf{RL instability.} Learning a policy on top of a video diffusion model is finicky
        (reward hacking, credit assignment over long streams, training cost). Fallback: if the
        learned pipeline is unstable, the training-free rule is a shippable result on its own,
        though a weaker CVPR contribution.
  \item \textbf{Reward reliability.} The online reward proxies (source preservation, edit
        alignment, temporal consistency) may be gameable or mis-calibrated; a policy that
        maximizes a bad proxy can degrade perceptual quality.
\end{itemize}

\section{Changelog}

\subsection*{2026-08-11 --- Initial idea seeded}
Established the core proposal: treat source-injection scheduling in a frozen causal streaming
video editor (StreamEdit setting) as a closed-loop control problem, exploiting the realized
causal state (decoded frames + KV cache) that a non-causal two-branch editor cannot access
single-pass. Chose a two-part structure: (1) a training-free online feedback rule as the
\emph{phenomenon / lower bound / ablation}, and (2) a \emph{learned injection-control policy}
as the headline method. Rationale for favoring the learned controller as the CVPR candidate:
the training-free-only version is vulnerable to the ``just a hand-tuned heuristic schedule''
objection and offers a thin mechanism, whereas the learned policy can make and test the strong
claim of approaching the oracle per-clip schedule. Recorded the oracle-schedule comparison as
the designated falsifier for the central claim, and flagged that causality must be load-bearing
in the controller's state or the contribution reduces to existing RL-over-guidance work
(``Guidance Is Not a Hyperparameter''). Seeded relation-to-prior-work from a verified
literature search (adaptive-guidance-control line; trained streaming editors LiveEdit /
JoyAI-Video-Edit / AnchorEdit; drift/exposure-bias line Self Forcing / TetherCache).
Open direction not yet committed: speculative/anytime ``draft-and-verify'' editing over the
injection knob, deferred as a possible second act.

\subsection*{2026-08-11 --- Sharpened falsifiers: oracle-per-clip bar and the GINH reduction}
Tightened three positioning points after a critique pass. (i) \emph{Corrected the falsifier for
the closed-loop claim}: beating dataset-averaged fixed-best is necessary but not sufficient, as
it only proves the schedule should be clip-dependent (capturable by cheap best-of-$K$ selection,
no causality). The proof of the \emph{within-clip} closed-loop gap lands only against
oracle-per-clip; updated claim 1, the method's subsection (1), and the oracle-gap open question
accordingly. (ii) \emph{Named GINH as the central threat explicitly}: the headline learned policy
\emph{is} GINH's abstraction (state-conditioned RL policy over a guidance-like scalar, sparse
terminal reward, frozen generator), so a reviewer can reduce it to ``GINH on a new knob in a new
modality.'' (iii) \emph{Reframed the causal-state defense}: ``we condition on causal state'' is
not an escape because GINH already conditions on state; our distinct contribution is the
\emph{efficiency} of cheap single-pass feedback, which upgrades to method-novelty only if the
causal state is shown \emph{necessary} (timestep-only $<$ state-conditioned). Elevated that
ablation to ``the paper'' in claim 3 and the relation-to-prior-work and open-questions sections.
No change to the method mechanism; these edits harden the novelty spine against the GINH
reduction.

\subsection*{2026-08-11 --- Expanded the two speculative flags in Method (2)}
Unpacked the reward-shaping and lookahead \textbf{(speculative)} flags inline with their pro/con
reasoning, so the design forks are legible without the transcript. Reward: terminal (sparse,
honest, hard credit assignment) vs.\ dense (shaped, easier, proxy/reward-hacking risk), default
sparse-terminal. Lookahead (MPC-lite): the try-$K$-candidates-and-pick idea, its place on the
fallback ladder (can be training-free), and the latency tension that makes it an open
experiment. Both forks remain undecided --- flags kept, no default committed beyond the reward
one.
\end{document}